% Part: first-order-logic
% Chapter: beyond
% Section: intuitionistic-logic

\documentclass[../../../include/open-logic-section]{subfiles}

\begin{document}

\olfileid{fol}{byd}{il}

\olsection{Intuïtionistische logica}

In tegenstelling tot tweede-ordelogica en hogere-ordelogica is
intuïtionistische eerste-ordelogica een beperking van de klassieke variant,
bedoeld om een meer ``constructieve'' redeneerwijze te modelleren. De volgende
voorbeelden kunnen enkele achterliggende motieven verduidelijken.

Stel dat iemand op een dag meedeelt een natuurlijk getal~$x$ te hebben bepaald
met de volgende eigenschap: als $x$ een priemgetal is, is de Riemann-hypothese
waar, en als $x$ samengesteld is, is de Riemann-hypothese onwaar. Dat lijkt
uitstekend nieuws. Of de Riemann-hypothese waar is, is een van de grote open
vragen in de wiskunde, en hier lijkt het probleem te zijn teruggebracht tot een
berekening, namelijk bepalen of één bepaald getal priem is.

Wat is de magische waarde van $x$? De beschrijving luidt: $x$ is het natuurlijke
getal dat gelijk is aan $7$ als de Riemann-hypothese waar is, en anders aan $9$.

Verontwaardigd eist men zijn geld terug. Klassiek gezien bepaalt deze
beschrijving inderdaad één unieke waarde van $x$; wat men werkelijk verlangt,
is echter een waarde van $x$ die \emph{expliciet} wordt gegeven.

Beschouw als een ander, wellicht minder gekunsteld voorbeeld de volgende vraag.
We weten dat een irrationaal getal tot een rationale macht kan worden verheven
met een rationaal resultaat. Zo is $\sqrt{2}^2 = 2$. Minder duidelijk is of een
irrationaal getal tot een \emph{irrationale} macht kan worden verheven met een
rationaal resultaat. De volgende stelling beantwoordt deze vraag bevestigend:

\begin{thm}
Er bestaan irrationale getallen $a$ en $b$ waarvoor $a^b$ rationaal is.
\end{thm}

\begin{proof}
Beschouw $\sqrt{2}^{\sqrt{2}}$. Als dit getal rationaal is, zijn we klaar: neem
$a = b = \sqrt{2}$. Anders is het irrationaal. Dan geldt
\[
(\sqrt{2}^{\sqrt{2}})^{\sqrt{2}} = \sqrt{2}^{\sqrt{2} \cdot
  \sqrt{2}} = \sqrt{2}^2 = 2,
\]
en dit is zeker rationaal. Neem in dit geval dus $a$ gelijk aan
$\sqrt{2}^{\sqrt{2}}$, en neem $b$ gelijk aan~$\sqrt 2$.
\end{proof}

Is dit een geldig bewijs? De meeste wiskundigen vinden van wel. Toch is er ook
hier iets onbevredigends: we hebben het bestaan bewezen van een paar reële
getallen met een bepaalde eigenschap, zonder te kunnen zeggen \emph{welk} paar
het is. Hetzelfde resultaat kan worden bewezen op een manier waarop het paar
$a$, $b$ \emph{wel} in het bewijs wordt gegeven: neem $a = \sqrt{3}$ en
$b = \log_3 4$. Dan
\[
a^b = \sqrt{3}^{\log_3 4} = 3^{1/2 \cdot \log_3 4} = (3^{\log_3
  4})^{1/2} = 4^{1/2}= 2,
\]
omdat $3^{\log_3 x} = x$.

Intuïtionistische logica is ontworpen om een redeneerwijze te modelleren waarin
stappen als die in het eerste bewijs niet zijn toegestaan. Het bestaan bewijzen
van een $x$ die aan~$!A(x)$ voldoet, betekent dat een concrete~$x$ moet worden
gegeven, samen met een bewijs dat deze aan $!A$ voldoet, zoals in het tweede
bewijs. Om te bewijzen dat $!A$ of $!B$ geldt, moet men de ene of de andere
formule kunnen bewijzen.

Formeel ontstaat intuïtionistische eerste-ordelogica door !!a{derivation} systeem
voor eerste-ordelogica op een bepaalde manier te beperken. Evenzo bestaan er
intuïtionistische varianten van tweede-ordelogica en hogere-ordelogica. Wiskundig
gezien zijn dit eenvoudig formele deductieve systemen, maar zij zijn, zoals
gezegd, bedoeld om een bepaalde vorm van wiskundig redeneren te modelleren. Men
kan deze opvatten als de redeneerwijze die door een bepaalde
wiskundefilosofische positie wordt gerechtvaardigd, zoals het intuïtionisme van
Brouwer; als een concretere en bevredigender vorm van wiskundig redeneren, in de
geest van het constructivisme van Bishop; en men kan betwisten of de formele
beschrijving de informele motivatie adequaat weergeeft. Ongeacht de ingenomen
filosofische positie kan intuïtionistische logica echter als een formeel
gepresenteerde logica worden bestudeerd, en om uiteenlopende redenen vinden veel
wiskundige logici dat interessant.

De intuïtionistische connectieven hebben een informele constructieve
interpretatie die gewoonlijk de BHK-interpretatie heet, naar Brouwer, Heyting en
Kolmogorov. Zij luidt als volgt: een bewijs van $!A \land !B$ bestaat uit een
bewijs van $!A$, gekoppeld aan een bewijs van $!B$; een bewijs van $!A \lor !B$
bestaat uit een bewijs van $!A$ of een bewijs van $!B$, waarbij expliciet bekend
is welk van beide het geval is; een bewijs van $!A \lif !B$ bestaat uit een
procedure die een bewijs van $!A$ omzet in een bewijs van~$!B$; een bewijs van
$\lforall[x][!A(x)]$ bestaat uit een procedure die een bewijs van $!A(x)$
oplevert voor elke waarde van~$x$; en een bewijs van $\lexists[x][!A(x)]$ bestaat uit
een waarde van~$x$, samen met een bewijs dat deze waarde aan~$!A$ voldoet. Deze
interpretatie kan computationeel worden beschreven met het zogeheten
``Curry--Howard-isomorfisme'' of het ``!!{formula}s-als-typenparadigma'': vat
!!a{formula} op als de specificatie van een bepaald gegevenstype, en bewijzen als
computationele objecten van zulke typen die laten zien dat de overeenkomstige
!!{formula} waar is.

Intuïtionistische logica wordt vaak opgevat als klassieke logica ``zonder'' het
principe van de uitgesloten derde. De volgende stelling preciseert dit.
\begin{thm}
Intuïtionistisch zijn de volgende axiomaschema's equivalent:
\begin{enumerate}
\item $(\lnot !A \lif \lfalse) \lif !A$.
\item $!A \lor \lnot !A$
\item $\lnot \lnot !A \lif !A$
\end{enumerate}
\end{thm}
Instanties van elk van deze schema's uit een van de andere afleiden is een goede
oefening in de intuïtionistische logica.

De eerste deductieve systemen voor de intuïtionistische propositielogica, bedoeld
als formaliseringen van Brouwers intuïtionisme, werden onafhankelijk van elkaar
door Kolmogorov, Glivenko en Heyting opgesteld. De eerste formalisering van de
intuïtionistische eerste-ordelogica, en van delen van de intuïtionistische
wiskunde, is van Heyting. Hoewel een aantal klassiek geldige schema's
intuïtionistisch niet geldig is, zijn vele dat wel.

De \emph{dubbele-negatievertaling} beschrijft een belangrijke relatie tussen
klassieke en intuïtionistische logica. Zij wordt als volgt inductief gedefinieerd
(vat $!A^N$ op als de ``intuïtionistische'' vertaling van de klassieke
!!{formula}~$!A$):
\begin{align*}
!A^N & \ident \lnot\lnot !A \quad \text{voor atomaire !!{formula}s $!A$} \\
(!A \land !B)^N & \ident (!A^N \land !B^N) \\
(!A \lor !B)^N & \ident  \lnot\lnot (!A^N \lor !B^N) \\
(!A \lif !B)^N & \ident (!A^N \lif !B^N) \\
(\lforall[x][!A])^N & \ident \lforall[x][!A^N] \\
(\lexists[x][!A])^N & \ident \lnot\lnot\lexists[x][!A^N]
\end{align*}
Kolmogorov en Glivenko gaven varianten van deze vertaling voor de
propositielogica; voor de predicatenlogica is zij onafhankelijk van elkaar door
G\"odel en Gentzen gegeven. Er geldt:

\begin{thm}
\begin{enumerate}
\item $!A \liff !A^N$ is klassiek bewijsbaar.
\item Als $!A$ klassiek bewijsbaar is, dan is $!A^N$ intuïtionistisch
  bewijsbaar.
\end{enumerate}
\end{thm}

Men kan zich nu de volgende dialoog voorstellen. Klassieke wiskundige: ``Ik heb
$!A$ bewezen!'' Intuïtionistische wiskundige: ``Uw bewijs is niet geldig. Wat u
werkelijk hebt bewezen, is $!A^N$.'' Klassieke wiskundige: ``Ook goed!'' Voor de
klassieke wiskundige maakt de intuïtionist een haarkloverij van de zaak, want de
twee zijn equivalent. De intuïtionist houdt echter vol dat er een verschil is.

Merk op dat de bovenstaande vertaling uitsluitend de zuivere logica betreft. Zij
beantwoordt niet de vraag welke \emph{niet-logische} axioma's voor de klassieke en
intuïtionistische wiskunde gepast zijn, noch hoe die axioma's zich tot elkaar
verhouden. De volgende geringe uitbreiding van de bovenstaande stelling geeft
echter enige nuttige informatie:

\begin{thm}
Als $\Gamma$ klassiek $!A$ bewijst, dan bewijst $\Gamma^N$
intuïtionistisch $!A^N$.
\end{thm}

Met andere woorden: als $!A$ klassiek bewijsbaar is uit bepaalde hypothesen, dan
is $!A^N$ bewijsbaar uit hun dubbele-negatievertalingen.

Om aan te tonen dat een zin of propositionele !!{formula} intuïtionistisch geldig
is, volstaat een bewijs. Maar hoe toont men aan dat zij niet geldig is? Daarvoor
is een correcte, en bij voorkeur volledige, semantiek nodig. De
Kripke-semantiek is daarvoor zeer geschikt.

We kunnen hetzelfde programma volgen als voor de klassieke logica: de semantiek
definiëren en correctheid en volledigheid bewijzen. Het is echter nuttig op het
volgende verschil te wijzen. Bij de klassieke logica was de semantiek in zekere
zin de ``vanzelfsprekende'', die impliciet in de betekenis van de connectieven
besloten ligt. Voor de Kripke-semantiek kan weliswaar een intuïtieve motivatie
worden gegeven, maar zij voelt niet op dezelfde manier onvermijdelijk aan.
Bovendien is het begrip klassieke !!{structure} een natuurlijk wiskundig begrip.
Men kan het begrip !!a{structure} daarom gebruiken als hulpmiddel om klassieke
eerste-ordelogica te bestuderen, of de klassieke eerste-ordelogica gebruiken om
wiskundige !!{structure}s te bestuderen. Kripke-!!{structure}s kunnen daarentegen
alleen als logische constructies worden beschouwd; zij lijken geen zelfstandig
wiskundig belang te hebben.

Een Kripke-!!{structure}~$\mModel M = \tuple{W, R, V}$ voor een propositionele
taal bestaat uit een verzameling~$W$, een partiële ordening~$R$ op~$W$ met een
kleinste !!{element}, en een ``monotone'' toekenning van propositionele
variabelen aan de !!{element}s van~$W$. De intuïtie is dat de !!{element}s van
$W$ ``werelden'' of ``kennistoestanden'' voorstellen; een element $v \geq u$
stelt een ``mogelijke toekomstige toestand'' van~$u$ voor; en de propositionele
variabelen die aan~$u$ zijn toegekend, zijn de proposities waarvan in toestand
~$u$ bekend is dat ze waar zijn. De forceringsrelatie $\mSat{M}{!A}[w]$ breidt
deze relatie vervolgens uit tot willekeurige !!{formula}s van de taal; lees
$\mSat{M}{!A}[w]$ als ``$!A$ is waar in toestand~$w$.'' De relatie wordt als
volgt inductief gedefinieerd:
\begin{enumerate}
\item $\mSat{M}{\Obj p_i}[w]$ desda $\Obj p_i$ een van de propositionele
  variabelen is die aan~$w$ zijn toegekend.
\item $\mSat/{M}{\lfalse}[w]$.
\item $\mSat{M}{(!A \land !B)}[w]$ desda $\mSat{M}{!A}[w]$ en $\mSat{M}{!B}[w]$.
\item $\mSat{M}{(!A \lor !B)}[w]$ desda $\mSat{M}{!A}[w]$ of $\mSat{M}{!B}[w]$.
\item $\mSat{M}{(!A \lif !B)}[w]$ desda voor alle $w' \geq w$ geldt: als
  $\mSat{M}{!A}[w']$, dan $\mSat{M}{!B}[w']$.
\end{enumerate}
Een goede oefening is aan te tonen dat $\lnot (p \land q) \lif
(\lnot p \lor \lnot q)$ niet intuïtionistisch geldig is, door een
Kripke-!!{structure} te construeren die een tegenvoorbeeld levert.

\end{document}
